\chapter{Package Framework}\label{ch:packageframework}

A central contribution of this thesis is the development of
\textit{AntCalc}\footnote{Source code available at
\url{https://github.com/HenFar/AntCalc/tree/release/0.3.x}.}, a \textit{Wolfram Language}
package created from scratch as part of this thesis work to automate the symbolic
construction and integration of QCD antenna functions. The package provides direct access to
both unintegrated and integrated antenna functions and implements the computational framework
described in the previous chapter. The results presented in this thesis were produced with
\textit{AntCalc}, version 0.3.0$\beta$2, corresponding to commit \texttt{4ceb991}.

\textit{AntCalc} makes use of \textit{FeynCalc} \cite{Feyncalc1,Feyncalc2,Feyncalc3,Feyncalc4}
for symbolic amplitude manipulation and interfaces with \textit{LiteRed2} \cite{Litered} for
IBP reduction.
Its internal architecture is organised as a
profile-driven route system: route-specific information is centralised in
profiles, which guide the selection and configuration of the appropriate build,
reduction, and integration workflows. This establishes a common contract
between the package components, reducing the amount of route-specific logic
that must be duplicated when existing functionality is extended or new routes
are introduced.

The package supports massless and massive antenna routes. All routes were validated against
the established literature, with all massless routes being validated against
\cite{Gehrmann-DeRidder:2005btv} except for the two-loop antenna set, which was validated
against \cite{Gehrmann-DeRidder:2004ttg}. The single massive route results were validated
against \cite{GehrmannDeRidder:2009fz}. All routes are public: none is held back as a 
partial or experimental implementation, and each produces stable and reproducible 
results notwithstanding the current beta version. The release-acceptance procedure 
that ensures this is described in Section~\ref{sec:validationDiagnostics}.

\textit{AntCalc}'s inner working structure can be described using the diagram in Fig.~\ref{fig:antcalc-architecture}.

\begin{figure}[h!]
\centering
\begin{tikzpicture}[
  font=\small,
  >=Stealth,
  stage/.style={
    draw=black!65,
    rounded corners=2pt,
    fill=white,
    align=center,
    minimum height=10mm,
    text width=19mm,
    inner sep=3pt
  },
  profile/.style={
    stage,
    fill=violet!7,
    text width=34mm,
    minimum height=12mm
  },
  public/.style={
    stage,
    fill=blue!7,
    text width=24mm
  },
  flow/.style={->, line width=.75pt},
  config/.style={->, dashed, gray!75, line width=.6pt}
]

% main lines
\draw[black, thick] (0,0) -- (10,0);
\draw[black, thick] (0,-5) -- (0,0);
\draw[black, thick] (5,0) -- (5,-5);
\draw[black, thick] (10,-5) -- (10,0);

% back and frontend backgrounds
\begin{pgfonlayer}{background}
  % Draws a light blue rectangle behind a specific region
  \fill[blue!10, rounded corners] (-2,1) rectangle (12,-4);
    
  % Or a light gray box behind another region
  \fill[black!5, rounded corners] (-2,-4) rectangle (12,-6);
\end{pgfonlayer}

% back and frontend labels
\node[label={[text=blue!60]above:\textit{Backend}}] at (11,-4) {};
\node[label={[text=gray]below:\textit{Frontend}}] at (11,-4) {};

% interceptions
\filldraw[black] (0,0) circle (2pt) node[anchor=south, yshift = 2pt]{Route Profile};
\filldraw[black] (5,0) circle (2pt) node[anchor=south]{BuildAntenna[]};
\filldraw[black] (10,0) circle (2pt) node[anchor=south]{IntegrateAntenna[]};

% input node
\node[draw, fill=black, isosceles triangle, isosceles triangle apex angle=60, inner sep=2pt, shape border rotate=+90, label=below:Input Key] at (0,-5) {};

% output nodes
\node[draw, fill=black, isosceles triangle, isosceles triangle apex angle=60, inner sep=2pt, shape border rotate=-90, label=below:Default Build Output] at (5,-5) {};
\node[draw, fill=black, isosceles triangle, isosceles triangle apex angle=60, inner sep=2pt, shape border rotate=-90, label=below:Default Integrate Output] at (10,-5) {};

% midway points
\node[draw, fill=black, isosceles triangle, isosceles triangle apex angle=60, inner sep=2pt, shape border rotate=+90, label={[label distance = -5pt, rotate=90, anchor=south]below:\textit{Resolving...}}] at (0,-2.5) {};
\node[draw, fill=black, isosceles triangle, isosceles triangle apex angle=60, inner sep=2pt, shape border rotate=-90, label={[label distance = -5pt, rotate=90, anchor=south]below:\textit{Formatting...}}] at (5,-2.5) {};
\node[draw, fill=black, isosceles triangle, isosceles triangle apex angle=60, inner sep=2pt, shape border rotate=-90, label={[label distance = -5pt, rotate=90, anchor=south]below:\textit{Expanding...}}] at (10,-2.5) {};
\node[draw, fill=black, isosceles triangle, isosceles triangle apex angle=60, inner sep=2pt, shape border rotate=-45] at (2.5,0) {};
\node[draw, fill=black, isosceles triangle, isosceles triangle apex angle=60, inner sep=2pt, shape border rotate=-45, label={[align=center]below:\textit{AntennaObject}\\\textit{(IntegrableForm -\textgreater \,True)}}] at (7.5,0) {};

\end{tikzpicture}
\caption{Profile-driven architecture of \textit{AntCalc}. Route profiles configure the build and integration workflows, while the build-side route context carries the evolving calculation data. The \texttt{AntennaObject} provides the integration-ready handoff between the two stages.}
\label{fig:antcalc-architecture}
\end{figure}


\section{Design Philosophy and Scope}\label{sec:philosophy}

\textit{AntCalc} was built to combine the symbolic amplitude capabilities of the 
FeynCalc environment with the IBP-reduction capabilities of LiteRed2, adapting and 
orchestrating the parts of those systems required for the construction and integration
of QCD antenna functions. 

The central design objective of the package is to return antenna functions in the
conventional forms required for higher-order QCD calculations using the antenna
subtraction method.
Though it supports a wide
variety of options to inspect and configure the backend operations, the default outputs
of its two central functions are designed to reflect this objective.
For example, the unintegrated one-loop three-parton antennae $\{A_3^1, \tilde A_3^1, \hat A_3^1\}$ returned by
\texttt{BuildAntenna[\ldots]} are expressed after Passarino-Veltman reduction, in terms of
scalar one-loop functions and the external kinematic invariants, following the conventional
representation used in the literature \cite{Gehrmann-DeRidder:2005btv}. By default, \texttt{IntegrateAntenna[\ldots]} returns
the integrated antennae as Laurent series in $\epsilon$, with the truncation order chosen to
match the conventions used in the antenna literature: the NLO antennae $\mathcal A_2^1$ and
$\mathcal A_3^0$ are expanded through $\mathcal O(\epsilon^2)$, while the NNLO antennae
$\mathcal A_4^0$, $\mathcal A_3^1$, $\mathcal A_2^2$, $\mathcal B_4^0$, and $\mathcal C_4^0$ are
returned through $\mathcal O(\epsilon^0)$ \cite{Gehrmann-DeRidder:2005btv,Gehrmann-DeRidder:2004ttg}.
Furthermore, the massive NLO tree-level antenna $\mathcal A_3^{0,(m_q)}$ is given up to
$\mathcal O(\epsilon)$ \cite{GehrmannDeRidder:2009fz}.

Its structure was also designed to be modular and extensible.
Each of the two central functions (\texttt{BuildAntenna} and \texttt{IntegrateAntenna}) is 
implemented as an encapsulated public module that communicates with the rest of the codebase
using the profile system: the route profile selects and configures the appropriate 
workflows and the route context carries the evolving results and diagnostics between 
stages.
Using route profiles and route contexts, as described in Section~\ref{sec:profile},
facilitates the addition of further modules while reducing duplicated logic and
direct dependencies. These new modules need not be
compatible with previous implemented functions through function-specific 
couplings. They need only conform with the profile system, being selected and 
configured by route profiles, and accepting inputs from and delivering outputs 
to the route contexts. 

Currently, \textit{AntCalc} is organised by mass scale, into massless and massive routes.
All massless quark-antiquark final-final antenna functions through NNLO have been made public.
Alongside these, a single massive route, $A_3^0$, is finalised, validated, and likewise public.
This is the first massive antenna of what is expected to be a complete
expansion of \textit{AntCalc}'s capabilities into the massive regime for quark-antiquark
final-final antennae.
The antenna functions currently implemented are summarised in Table~\ref{tab:antennaTiers}.

\begin{table}[h!]
\centering
\renewcommand{\arraystretch}{1.3}
  \begin{tabular}{|c|c|c|}
   \hline
   \textbf{Perturbative Order} & \textbf{Antenna} & \textbf{Validated Against} \\
   \hline
   LO & $A_2^0$ & Born normalisation \\ \hline
   NLO & $A_2^1$, $A_3^0$ & \cite{Gehrmann-DeRidder:2005btv,Gehrmann-DeRidder:2004ttg} \\ \hline
   NNLO & $A_3^1$, $\tilde A_3^1$,
   $\hat A_3^1$, $A_4^0$, $\tilde A_4^0$, $B_4^0$, $C_4^0$ & \cite{Gehrmann-DeRidder:2005btv,Gehrmann-DeRidder:2004ttg} \\ \hline
   NNLO & $A_2^2$, $\tilde A_2^2$, $\hat A_2^2$, $\breve A_2^2$ & \cite{Gehrmann-DeRidder:2004ttg} \\ \hline
   NLO, massive & $A_3^{0,(m_q)}$ & \cite{GehrmannDeRidder:2009fz} \\ \hline
  \end{tabular}
  \caption{Overview of the antenna functions currently implemented in \textit{AntCalc}, by
  perturbative order and validation source.}
\label{tab:antennaTiers}
\end{table}

\section{Profile-Driven Architecture}\label{sec:profile}

As described in Section~\ref{sec:philosophy}, \textit{AntCalc} relies on a 
profile-driven routing system. This system maps the route key to its route-specific 
profile and workflow. 

A public antenna route key is a three-element list which describes a canonical route
whose execution may be refined by options. The route key takes the form,
\begin{equation}
    \{\text{Antenna Type}, \text{Final-State Multiplicity}, \text{Loop Order}\}.
\end{equation}
As an example, the route key for
the one-loop three-parton
antenna set associated with the full-process amplitude $M_{qg\bar q}^1$ is $\{\text A,3,1\}$. This key
represents the complete set of colour contributions, comprising the leading-colour
antenna $A_3^1$, the subleading-colour antenna $\tilde A_3^1$, and the quark-loop
contribution $\hat A_3^1$. The default public output is therefore the ordered list
$\{A_3^1, \tilde A_3^1, \hat A_3^1\}$.

The package links this key to the profile it identifies and follows that profile's 
directions on how the route should be constructed and/or integrated. The data produced 
while the route executes is also stored in the route context association. 
The profile system consists of a main profile and two further profiles: one build-side,
handling reduction, and one integration-side.
The principal build-side profile is
\texttt{AntennaProfile[key]}, which specifies information such as antenna type, 
multiplicity and loop order, colour normalisation, extraction rules, components, etc.
\texttt{AntennaReductionProfile[key]} is separate, but still build-side. It specifies 
the Passarino-Veltman-reduction preferences for the appropriate routes. In this case, up to
NNLO, these are the $\{\text A,2,1\}$ and 
$\{\text A,3,1\}$ loop routes. \\ \texttt{AntennaIntegrationProfile[key]} is the 
integration-side profile, which specifies its side's backend, IBP basis 
families, expansion depth, etc. The specifics of each stage's profiles and options 
will be discussed in their respective Sections~\ref{sec:buildStage} and 
\ref{sec:integrateStage}.

While profiles select and configure computation paths, the workflows themselves 
perform the calculations. The route dispatcher and workflow code consult the profiles 
to select and configure each stage. As a representative example, which will be further 
explained and explored in the two upcoming sections, let us consider the call 
\texttt{BuildAntenna[A,2,1]}. Its full route profile is,

\begin{mdframed}[backgroundcolor=gray!10, linewidth=0pt]
\begin{lstlisting}
AntennaProfile[{A, 2, 1}] :=
    <|
        "Key" -> {A, 2, 1},
        "Name" -> "A21",
        "AntennaType" -> A,
        "NumFinalParticles" -> 2,
        "LoopOrder" -> 1,
        "TreeAmplitude" -> AntennaAmplitude[{A, 2, 0}],
        "BornInterference" -> BornInterference[],
        "Extraction" -> "LoopScalar",
        "ColourNorm" -> colourNorm,
        "ExcludedParticles" -> 
            {S[_], V[1], V[2], V[3], U[1], U[2], U[3], U[4]},
        "DropSumOver" -> True,
        "ReductionProfile" -> AntennaReductionProfile[{A, 2, 1}],
        "ConventionProfile" -> BuildAntennaConventionProfile[{A, 2, 1}]
    |>
\end{lstlisting}
\end{mdframed}
In this profile, the key, name, antenna type, number of final particles (multiplicity),
and loop order have been discussed before. The settings used in this route are explained 
in Table~\ref{tab:A21Sets}.

\begin{table}[h]
    \centering
    \renewcommand{\arraystretch}{1.3}
    \begin{tabular}{|Sc|Sc|}
        \hline
        \textbf{Item} & \textbf{Use}\\
        \hline
        \texttt{TreeAmplitude} & Selects the tree-level amplitude. In this case, it is $\mathcal M_2^0$.\\
        \hline
        \texttt{BornInterference} & Generates the general normalisation as $\left|\mathcal M_2^0\right|^2$. \\ 
        \hline
        \texttt{Extraction} & Selects how the components are extracted. \\
        \hline
        \texttt{ColourNorm} & Colour normalisation: $N-1/N$. \\
        \hline
        \texttt{ExcludedParticles} & Restricts \textit{FeynArts} internal particle insertions. \\
        \hline
        \texttt{DropSumOver} & Replaces every \texttt{SumOver[]} statement with $1$. \\
        \hline
        \texttt{ReductionProfile} & Points to the reduction profile. \\
        \hline
        \texttt{ConventionProfile} & Points to the normalisation-convention ruleset profile. \\
        \hline
    \end{tabular}
    \caption{\centering Overview of the settings for the $A_2^1$ route used as a representative example of the route system used in \textit{AntCalc}.}
    \label{tab:A21Sets}
\end{table}

\section{The Build Stage}\label{sec:buildStage}

The build stage constructs the unintegrated antenna functions from their corresponding matrix
elements. Its main public interface is
\texttt{BuildAntenna[type, multiplicity, loop order]}. The workflow begins with a route key
specified through a public call, for example \texttt{BuildAntenna[A, 3, 1]}. From this key,
\textit{AntCalc} identifies the appropriate construction workflow and the ingredients required
for the corresponding antenna set. The relevant amplitudes are then constructed and combined
according to the route definition, after which the resulting expressions are processed into the
required unintegrated antenna functions. Internally, the intermediate results are stored in an
association containing the different components of the calculation. The final step reformats
this internal representation into the antenna expression, or ordered set of antenna components,
returned to the user. The build workflow can therefore be summarised as follows,

\begin{mdframed}[backgroundcolor=gray!10, linewidth=0pt]
\begin{lstlisting}
BuildAntenna[type, multiplicity, loop order]
    -> route key;
    -> loop-order dispatcher;
    -> route-specific build workflow;
    -> build-data association / route context;
    -> public-output boundary;
    -> output: expression, component list, AntennaObject, diagnostics, 
       or run record.
\end{lstlisting}
\end{mdframed}

The loop-order dispatcher is a central element of this workflow. Using the route key, it
directs each antenna to the appropriate computational path, since different loop orders
require different computational treatments, most notably the tensor reduction required for
the loop-momentum dependence of one- and two-loop antennae.

Functionally, the route context for this stage holds the profile, the generated amplitude, the
computed and normalised interference, the extracted components, and the diagnostics; some routes
retain additional route-specific intermediates.

A diagram in the spirit of the one in Fig.~\ref{fig:antcalc-architecture}, but only relative to 
the build stage is shown below, in Fig.~\ref{fig:antcalc-architecture-buildstage}.

\begin{figure}[h!]
\centering
\begin{tikzpicture}[
  font=\small,
  >=Stealth,
  stage/.style={
    draw=black!65,
    rounded corners=2pt,
    fill=white,
    align=center,
    minimum height=10mm,
    text width=19mm,
    inner sep=3pt
  },
  profile/.style={
    stage,
    fill=violet!7,
    text width=34mm,
    minimum height=12mm
  },
  public/.style={
    stage,
    fill=blue!7,
    text width=24mm
  },
  flow/.style={->, line width=.75pt},
  config/.style={->, dashed, gray!75, line width=.6pt}
]

% main lines
\draw[black, thick] (0,0) -- (7.5,0);
\draw[black, dashed] (7.5,0) -- (10,0);
\draw[black, thick] (0,-5) -- (0,0);
\draw[black, thick] (7.5,0) -- (7.5,-5);

% back and frontend backgrounds
\begin{pgfonlayer}{background}
  % Draws a light blue rectangle behind a specific region
  \fill[blue!10, rounded corners] (-2,1) rectangle (12,-4);
    
  % Or a light gray box behind another region
  \fill[black!5, rounded corners] (-2,-4) rectangle (12,-6);
\end{pgfonlayer}

% back and frontend labels
\node[label={[text=blue!60]above:\textit{Backend}}] at (11,-4) {};
\node[label={[text=gray]below:\textit{Frontend}}] at (11,-4) {};

% interceptions
\filldraw[black] (0,0) circle (2pt) node[anchor=south, yshift = 2pt]{Route Profile};
\filldraw[black] (7.5,0) circle (2pt) node[anchor=south]{BuildAntenna[]};

% input node
\node[draw, fill=black, isosceles triangle, isosceles triangle apex angle=60, inner sep=2pt, shape border rotate=+90, label=below:Input Key] at (0,-5) {};

% output nodes
\node[draw, fill=black, isosceles triangle, isosceles triangle apex angle=60, inner sep=2pt, shape border rotate=-90, label=below:Default Build Output] at (7.5,-5) {};

% midway points
\node[draw, fill=black, isosceles triangle, isosceles triangle apex angle=60, inner sep=2pt, shape border rotate=+90, label={[label distance = -5pt, rotate=90, anchor=south]below:\textit{Resolving...}}] at (0,-2.5) {};
\node[draw, fill=black, isosceles triangle, isosceles triangle apex angle=60, inner sep=2pt, shape border rotate=-90, label={[label distance = -5pt, rotate=90, anchor=south]below:\textit{Formatting...}}] at (7.5,-2.5) {};
\node[draw, fill=black, isosceles triangle, isosceles triangle apex angle=60, inner sep=2pt, shape border rotate=-45] at (3.75,0) {};
\node[draw, fill=black, isosceles triangle, isosceles triangle apex angle=60, inner sep=2pt, shape border rotate=-45, label={[align=center]below:\textit{AntennaObject}\\\textit{(IntegrableForm -\textgreater \,True)}}] at (10,0) {};

\end{tikzpicture}
\caption{Profile-driven architecture of \textit{AntCalc}'s Build stage.}
\label{fig:antcalc-architecture-buildstage}
\end{figure}


\subsection{Tree-Level Routes}\label{subsec:treeLevelRoutes}

As mentioned, the workflow handles distinct loop-order routes differently, though the differences only appear 
after the dispatcher is called. Tree-level routes like those relative to the $A_3^0$ antenna, start by retrieving 
the profile and computing the relevant 
process amplitude. It uses the profile's \texttt{Production} field to select the physical construction path,
and the \texttt{Extraction} and \texttt{ColourNorm} fields to select the extraction mode and 
the colour normalisation used for each antenna. The colour normalisation used by default 
has been described previously in Eq.~\ref{eq:colourNorm}, and the 
\texttt{Extraction} mode 
also introduces the Born normalisation.
The routes selected in the 
\texttt{Production} and \texttt{Extraction} fields
depend on the methodology required to compute the antenna functions and retrieve the components, if required. 
The public tree-level antenna functions are constructed using different production
and extraction modes, as described in 
Table~\ref{tab:treeLevelProdModes}.

\begin{table}[h!]
\centering
\renewcommand{\arraystretch}{1.3}
\resizebox{\textwidth}{!}{
  \begin{tabular}{|c|c|c|} 
   \hline
   \textbf{Production Mode} & \textbf{Extraction Mode} & \textbf{Current Route} \\ 
   \hline
   \texttt{SelfInterference} & \texttt{BornScalar} & $A_2^0$ \\ \hline
   \texttt{SelfInterference} & \texttt{TreeColourCoefficients} & $A_3^0$ \\ \hline
   \texttt{ColourOrderedAntenna} & \texttt{TreeColourCoefficients} & $A_4^0$, $\tilde A_4^0$ \\ \hline
   \texttt{SectorSelfInterference} & \texttt{ColourNormScalar} & $B_4^0$ \\ \hline
   \texttt{SectorSymmetrisedInterference} & \texttt{MinusTwoColourNormOverSUNN} & $C_4^0$ \\ 
   \hline
   \texttt{SelfInterference} & \texttt{TreeColourCoefficients} & $A_3^{0,(m_q)}$ \\ 
   \hline
  \end{tabular}
} 
  \caption{Overview of the tree-level production and extraction modes.}
\label{tab:treeLevelProdModes}
\end{table}

These production paths all lead to the extraction of the required antenna functions, but the way in which
they do it changes with the antenna's multiplicity. Each production and extraction mode is
\textit{AntCalc}'s operational implementation of the construction derived, in mathematical terms, for the
corresponding \textit{Current Route} entry in Subsection~\ref{subsec:antFamilies}. On the extraction side,
$A_2^0$ is set to $1$; both $A_3^0$ and the $A_4^0$-set antennae are extracted from the usual
$NX_n^l + \frac1N\tilde X_n^l$ structure; and the $B_4^0$ and $C_4^0$ antennae are extracted from the more
involved inner structure of the $M_{4q}^0$ full-process amplitude, requiring a correspondingly more complex
extraction process.

\subsection{One-Loop Routes}\label{subsec:oneLoopRoutes}

One-loop routes are dispatched to \texttt{BuildLoopAntennaData} where the primary backend
reduction method is selected through the nested \texttt{ReductionProfile}, which by default is
Passarino-Veltman reduction, as described in Section~\ref{subsec:PaVe}.

The production method consists of generating the tree-level and one-loop amplitudes, using \\
\texttt{TreeAmplitude} and \texttt{MAmpOneLoop}, respectively, and then forming the tree-loop 
interference. The interference is colour- and one-loop-normalised, as described in Eqs.~\ref{eq:loopNorm} and \ref{eq:colourNorm}.
This normalised interference is then reduced and its components extracted.
The production and extraction modes implemented for the supported one-loop routes
are summarised in Table~\ref{tab:oneLoopProdModes}, below.

\begin{table}[h!]
\centering
\renewcommand{\arraystretch}{1.3}
  \begin{tabular}{|c|c|c|} 
   \hline
   \textbf{Production Mode} & \textbf{Extraction Mode} & \textbf{Current Route} \\ 
   \hline
   \texttt{TreeAmplitude + MAmpOneLoop} & \texttt{LoopScalar} & $A_2^1$ \\ \hline
   \texttt{TreeAmplitude + MAmpOneLoop} & \texttt{LoopColourCoefficients} & $A_3^1$, $\tilde A_3^1$, $\hat A_3^1$ \\ \hline
  \end{tabular}
  
  \caption{Overview of the one-loop production and extraction modes.}
\label{tab:oneLoopProdModes}
\end{table}

The two one-loop routes make use of a hybrid production mode that combines the \texttt{TreeAmplitude} and
\texttt{MAmpOneLoop} objects. As shown in Table~\ref{tab:A21Sets}, the former is $\mathcal M_n^0$, while the 
latter is its one-loop counterparts $\mathcal M_n^1$. These two amplitudes are used as described in 
Subsection~\ref{subsec:antFamilies}. Their extraction modes are the one-loop analogous to the methods 
used for $A_2^0$ and $A_3^0$, as discussed in Subsection~\ref{subsec:treeLevelRoutes}. 

\subsection{Two-Loop Routes}

Loop-order-two routes are dispatched to \texttt{BuildTwoLoopAntennaData}. There is only, formally, 
one set of this kind up to NNLO, $A_2^2$. It is remarkably different from the rest, however. The four-component set is not generated like 
the $A_4^0$ and $A_3^1$ sets were, rather its four components come from two distinct source contributions, as discussed in 
Subsection~\ref{subsec:antFamilies}.
Their production and extraction modes are similar and handled as a bundle inside \textit{AntCalc},
but, as shown in Table~\ref{tab:twoLoopProdModes}, their sources differ,

\begin{table}[h!]
\centering
\renewcommand{\arraystretch}{1.3}
  \begin{tabular}{|c|c|c|} 
   \hline
   \textbf{Production Mode} & \textbf{Extraction Mode} & \textbf{Current Route} \\ 
   \hline
   \texttt{TwoLoopInterf} & \texttt{TwoLoopTTermComponents} & $A_2^2$, $\tilde A_2^2$, $\hat A_2^2$ \\ \hline
   \texttt{OneLoopSelf} & \texttt{TwoLoopTTermComponents} & $\breve A_2^2$ \\ \hline
  \end{tabular}
  
  \caption{Overview of the two-loop contributions and extraction modes.}
\label{tab:twoLoopProdModes}
\end{table}

One of the ways in which this route is particular inside \textit{AntCalc}'s architecture 
is illustrated by their different production, but equal extraction modes. 
This route may be schematically represented as shown in Fig.~\ref{fig:a22-build}, below.

\begin{figure}[h!]
\centering
\begin{tikzpicture}[
  font=\small,
  >=Stealth,
  stage/.style={
    draw=black!65,
    rounded corners=2pt,
    fill=white,
    align=center,
    minimum height=10mm,
    text width=19mm,
    inner sep=3pt
  },
  profile/.style={
    stage,
    fill=violet!7,
    text width=34mm,
    minimum height=12mm
  },
  public/.style={
    stage,
    fill=blue!7,
    text width=24mm
  },
  flow/.style={->, line width=.75pt},
  config/.style={->, dashed, gray!75, line width=.6pt}
]

% centre line
\draw[black, thick] (-0.5,-1) -- (1,-1);

\draw[black, thick] (10,-1) -- (11.5,-1);

% build-only lines
\draw[black, thick] (3,-2) -- (6.5,-2);
\draw[black, thick] (7.5,-2) -- (8,-2);

% build-only nodes

% integrate-only lines
\draw[black, thick] (3,0) -- (5,0);

\draw[black, thick] (6,0.8) -- (6.5,0.8);
\draw[black, thick] (6,0) -- (6.5,0);
\draw[black, thick] (6,-0.8) -- (6.5,-0.8);

\draw[black, thick] (7.5,0.8) -- (8,0.8);
\draw[black, thick] (7.5,0) -- (8,0);
\draw[black, thick] (7.5,-0.8) -- (8,-0.8);

% integrate-only nodes
\filldraw[black] (7,0.8) node[anchor=center]{$A_2^2$};
\filldraw[black] (7,0) node[anchor=center]{$\tilde{A}_2^2$};
\filldraw[black] (7,-0.8) node[anchor=center]{$\hat{A}_2^2$};
\filldraw[black] (7,-2) node[anchor=center]{$\breve{A}_2^2$};

% curvy lines
\draw (1,-1) to[out=0,in=180] (3,-2);
\draw (1,-1) to[out=0,in=180] (3,0);

\draw (5,0) to[out=0,in=180] (6,0.8);
\draw (5,0) to[out=0,in=180] (6,0);
\draw (5,0) to[out=0,in=180] (6,-0.8);

\draw (8,0.8) to[out=0,in=180] (10,-1);
\draw (8,0) to[out=0,in=180] (10,-1);
\draw (8,-0.8) to[out=0,in=180] (10,-1);
\draw (8,-2) to[out=0,in=180] (10,-1);

% back and frontend backgrounds
\begin{pgfonlayer}{background}
  % Draws a light blue rectangle behind a specific region
  \fill[black!10, rounded corners] (-2.5,2) rectangle (13.5,-3.5);
\end{pgfonlayer}

% diagram labels
% \node[label={[text=gray]right:\textit{PaVe}}] at (-2,-4) {};

% circle points
\filldraw[black] (-0.5,-1) circle (2pt) node[align=center, anchor=south, yshift = 2pt]{\texttt{BuildAntenna} \\ \texttt{[A,2,2]}};

\filldraw[black] (11.5,-1) circle (2pt) node[align=center, anchor=south, yshift = 2pt]{$\{A_2^2, \tilde{A}_2^2, \hat{A}_2^2, \breve{A}_2^2\}$};

% arrows
\node[draw, fill=black, circle, inner sep=2pt, label={[align=center]above:$\mathcal J_2^{2,[2\times 0]}$}] at (4,0) {};
\node[draw, fill=black, regular polygon, regular polygon sides=5, inner sep=2pt, label={[align=center]below:$\mathcal J_2^{2,[1\times 1]}$}] at (4,-2) {};

\node[label={[text=blue!80,align=center]above:$(N)$}] at (6.1,0.7) {};
\node[label={[text=blue!80,align=center]above:$(1/N)$}] at (6.1,-0.1) {};
\node[label={[text=blue!80,align=center]above:$(N_f)$}] at (6.1,-0.9) {};

\end{tikzpicture}
\caption{\centering Diagram showing the build-side workflow leading to the generation of the $A_2^2$ antenna set. The two subsets are separated by the interferences used in their generation.}
\label{fig:a22-build}
\end{figure}


The programme reads the profile and generates the two necessary amplitudes, for the
two separate loops and for the one-loop self-interference case, $\mathcal M_2^{2,[2\times 0]}$
and $\mathcal M_2^{2,[1\times 1]}$. The former amplitude generates three antenna-set 
components, while the latter produces only one, but all four are extracted using the 
same usual method used for tree-level and one-loop antenna functions, applied 
to two loops. 

\section{From the Build Stage to Integration}

The build stage described above produces the unintegrated antenna functions in terms of the
external kinematic invariants $s_{ij}$. This is the natural public representation of the
antenna and is suitable for its use as a local subtraction term in a numerical phase-space
integration, as discussed in Chapter~\ref{ch:physicsbackground}. A standard
\texttt{BuildAntenna[...]} call therefore returns either a single antenna expression or, where
several antenna-set components are present, an ordered list made up of those components.

For the analytic calculation of the integrated antennae, however, the final expression in
terms of Mandelstam invariants is not by itself sufficient. The integration stage requires
additional information from the construction of the antenna, including the route definition
and the integral structures needed for the subsequent reduction and phase-space integration.
\textit{AntCalc} therefore provides the option \texttt{IntegrableForm -> True}. With this
option, \texttt{BuildAntenna[...]} returns an \texttt{AntennaObject} that retains both the
constructed antenna and the additional information required by \texttt{IntegrateAntenna[...]}.
This object provides the interface between the build and integration stages.

At this point, \textit{AntCalc} can therefore provide the unintegrated antennae in the
kinematic form required for local subtraction. The next stage addresses the complementary part
of the antenna subtraction procedure: their analytic integration over the unresolved phase
space in dimensional regularisation. These integrated antennae are required to combine the
subtraction terms with the virtual contributions and achieve the cancellation of explicit
infrared poles. The implementation of this integration stage is described in the following
section.

\section{The Integration Stage}\label{sec:integrateStage}

The integration stage performs the analytic integration of the antenna functions constructed
in the build stage over the corresponding unresolved phase space in $d=4-2\epsilon$ dimensions,
as in Eq.~\ref{eq:antennaIntegral}. The integrated antennae are denoted with calligraphic letters.

Within \textit{AntCalc}, this stage is initiated by passing the \texttt{AntennaObject} produced
by \\ \texttt{BuildAntenna[..., IntegrableForm -> True]} to \texttt{IntegrateAntenna[...]}. A 
diagram representing this architecture and how it splits from an unchanged \texttt{BuildAntenna[...]} 
call is shown in Figure~\ref{fig:component-extraction}.

\begin{figure}[h!]
\centering
\begin{tikzpicture}[
  font=\small,
  >=Stealth,
  stage/.style={
    draw=black!65,
    rounded corners=2pt,
    fill=white,
    align=center,
    minimum height=10mm,
    text width=19mm,
    inner sep=3pt
  },
  profile/.style={
    stage,
    fill=violet!7,
    text width=34mm,
    minimum height=12mm
  },
  public/.style={
    stage,
    fill=blue!7,
    text width=24mm
  },
  flow/.style={->, line width=.75pt},
  config/.style={->, dashed, gray!75, line width=.6pt}
]

% centre line
\draw[black, thick] (0,-1) -- (2,-1);

% build-only lines
\draw[black, thick] (4,-2) -- (5,-2);
\draw[black, dashed] (5,-2) -- (6,-2);

\draw[black, dashed] (7, -1.5) -- (7.5,-1.5);
\draw[black, dashed] (7, -2) -- (7.5,-2);
\draw[black, dashed] (7, -2.5) -- (7.5,-2.5);

% build-only nodes
\filldraw[black] (8,-1.5) node[anchor=center]{$A_3^1$};
\filldraw[black] (8,-2) node[anchor=center]{$\tilde A_3^1$};
\filldraw[black] (8,-2.5) node[anchor=center]{$\hat A_3^1$};

% integrate-only lines
\draw[black, thick] (4,0) -- (6,0);

\draw[black, thick] (7,0.5) -- (7.5,0.5);
\draw[black, thick] (7,0) -- (7.5,0);
\draw[black, thick] (7,-0.5) -- (7.5,-0.5);

\draw[black, dashed] (8.5,0.5) -- (9,0.5);
\draw[black, dashed] (8.5,0) -- (9,0);
\draw[black, dashed] (8.5,-0.5) -- (9,-0.5);

\draw[black, dashed] (10,0) -- (10.5,0)
                    arc [start angle=90, end angle=0, radius=0.5]
                           -- (11,-2);

% integrate-only nodes
\filldraw[black] (8,0.5) node[anchor=center]{$\mathcal A_3^1$};
\filldraw[black] (8,0) node[anchor=center]{$\tilde{\mathcal A}_3^1$};
\filldraw[black] (8,-0.5) node[anchor=center]{$\hat{\mathcal A}_3^1$};

% curvy lines
\draw (2,-1) to[out=0,in=180] (4,-2);
\draw (2,-1) to[out=0,in=180] (4,0);

\draw (6,0) to[out=0,in=180] (7,0.5);
\draw (6,0) to[out=0,in=180] (7,0);
\draw (6,0) to[out=0,in=180] (7,-0.5);

\draw[black, dashed] (6,-2) to[out=0,in=180] (7, -1.5);
\draw[black, dashed] (6,-2) to[out=0,in=180] (7, -2);
\draw[black, dashed] (6,-2) to[out=0,in=180] (7, -2.5);

\draw[black, dashed] (9,0.5) to[out=0,in=180] (10, 0);
\draw[black, dashed] (9,0) to[out=0,in=180] (10, 0);
\draw[black, dashed] (9,-0.5) to[out=0,in=180] (10, 0);

% back and frontend backgrounds
\begin{pgfonlayer}{background}
  % Draws a light blue rectangle behind a specific region
  \fill[black!10, rounded corners] (-2,1.5) rectangle (14,-4);
\end{pgfonlayer}

% diagram labels
% \node[label={[text=gray]right:\textit{PaVe}}] at (-2,-4) {};

% circle points
\filldraw[black] (0,-1) circle (2pt) node[anchor=south, yshift = 2pt]{BuildAntenna[A,3,1]};
\filldraw[black] (11,-2) circle (2pt) node[align=center, anchor=north, yshift = -2pt]{$\{\mathcal A_3^1, \tilde{\mathcal A}_3^1, \hat{\mathcal A}_3^1\}$};
\filldraw[black] (5,-2) circle (2pt) node[anchor=north, yshift = -2pt]{$\{A_3^1, \tilde A_3^1, \hat A_3^1\}$};

% arrows
\node[draw, fill=black, isosceles triangle, isosceles triangle apex angle=60, inner sep=2pt, shape border rotate=-45, label={[align=center]above:\textit{AntennaObject}}] at (5,0) {};
\node[draw, isosceles triangle, isosceles triangle apex angle=60, inner sep=2pt, shape border rotate=-90] at (11,-1) {};

\end{tikzpicture}
\caption{\centering Diagram showing the component extraction steps and paths for the build-only and build-and-integrate routes. The combined route is shown above and the build-only route below it. The dashed lines represent the options to either separate the ordered lists into single components on the build-only route, or to merge the components into a single ordered list for the combined route. The $A_3^1$ antenna set is used as an example.}
\label{fig:component-extraction}
\end{figure}


The route key stored in this object determines the appropriate integration workflow, including the
required reduction procedure, master integrals, and normalisation conventions. After these
steps have been applied, \texttt{IntegrateAntenna[...]} returns the integrated antenna, or the
corresponding set of antenna components, expanded as a Laurent series in $\epsilon$. The
integration workflow is structured as follows,

\begin{mdframed}[backgroundcolor=gray!10, linewidth=0pt]
\begin{lstlisting}
IntegrateAntenna[AntennaObject]
    -> recover route key and resolve profile;
    -> choose integration backend;
    -> complete IBP/Passarino-Veltman reduction;
    -> substitute master integrals when applicable;
    -> extract public integrated antennae;
    -> output: selected component, list, diagnostics, run record, or 
       master-combination view.
\end{lstlisting}
\end{mdframed}

It is relevant to note that \texttt{IntegrateAntenna} requires the profile to complete its 
function like \texttt{BuildAntenna}, since it also encodes the relevant routing information 
for the reduction and integration process. 

Much like for \texttt{BuildAntenna}, the route state is not held in a single context object
but spread across several objects, accessed as required: the key and profile, the input
(unintegrated) antenna and selected component, the IBP bases and master integrals, the series
result, and the diagnostics.

An integration-stage-specific diagram in the spirit of the one in
Fig.~\ref{fig:antcalc-architecture} can also be drawn in further detail. 
That diagram is shown in Fig.~\ref{fig:antcalc-architecture-intstage}.

\begin{figure}[h!]
\centering
\begin{tikzpicture}[
  font=\small,
  >=Stealth,
  stage/.style={
    draw=black!65,
    rounded corners=2pt,
    fill=white,
    align=center,
    minimum height=10mm,
    text width=19mm,
    inner sep=3pt
  },
  profile/.style={
    stage,
    fill=violet!7,
    text width=34mm,
    minimum height=12mm
  },
  public/.style={
    stage,
    fill=blue!7,
    text width=24mm
  },
  flow/.style={->, line width=.75pt},
  config/.style={->, dashed, gray!75, line width=.6pt}
]

% main lines
\draw[black, thick] (0,0) -- (11,0);
\draw[black, thick] (11,0) -- (11,-10);

% back and frontend backgrounds
\begin{pgfonlayer}{background}
  % Draws a light blue rectangle behind a specific region
  \fill[blue!10, rounded corners] (-2.5,1) rectangle (14,-8);
    
  % Or a light gray box behind another region
  \fill[black!5, rounded corners] (-2.5,-8) rectangle (14,-12);

  % Red background for options
  \fill[red!15, rounded corners] (2.5,-1) rectangle (10,-11);
\end{pgfonlayer}

% back and frontend labels
\node[label={[text=blue!60]above:\textit{Backend}}] at (13,-8) {};
\node[label={[text=gray]below:\textit{Frontend}}] at (13,-8) {};
\node[label={[text=red!60]below:\textit{Options}}] at (3.3,-10.2) {};

% interceptions
\filldraw[black] (0,0) circle (2pt) node[align=center, anchor=north, yshift = -2pt]{AntennaObject \\ \textit{(IntegrateAntenna-\textgreater True)}};
\filldraw[black] (11,0) circle (2pt) node[anchor=south]{IntegrateAntenna[]};

% input node

% output nodes
\node[draw, fill=black, isosceles triangle, isosceles triangle apex angle=60, inner sep=2pt, shape border rotate=-90, label={[align=center, anchor=west, label distance=3pt]below:Default \\ Integrate \\ Output}] at (11,-10) {};

% midway points
\node[draw, fill=black, isosceles triangle, isosceles triangle apex angle=60, inner sep=2pt, shape border rotate=-90, label={[label distance = -5pt, rotate=90, anchor=south]below:\textit{Expanding...}}] at (11,-5) {};
\node[draw, fill=black, isosceles triangle, isosceles triangle apex angle=60, inner sep=2pt, shape border rotate=-45] at (5.5,0) {};

% Options
\draw[black, dashed] (2.1,0) -- (2.1,-1.6);
\draw[black, dashed] (2.1,-1.6) -- (2.85,-1.6);

\node[draw, fill=black, isosceles triangle, isosceles triangle apex angle=60, inner sep=2pt, shape border rotate=-45] at (2.85,-1.6) {};

\node[align=center, anchor=west] at (3.05,-1.7) {Convention \\ \& Expansion:};
\draw[black!80, thick] (3.8, -2.2) -- (3.8, -4);
\node[anchor=west] at (3.9,-2.45) {\texttt{ExpansionOrder}};
\node[anchor=west] at (3.9,-2.85) {\texttt{KinematicScale}};
\node[anchor=west] at (3.9,-3.25) {\texttt{NormalizeKinematicScale}};
\node[anchor=west] at (3.9,-3.65) {\texttt{ApplyFeynCalcMS}};

\node[align=center, anchor=west] at (3.05,-4.6) {IBP Basis \\ Control:};
\draw[black!80, thick] (3.8, -5.1) -- (3.8, -6.45);
\node[anchor=west] at (3.9,-5.35) {\texttt{BasisFamily}};
\node[anchor=west] at (3.9,-5.75) {\texttt{BasisRoot}};
\node[anchor=west] at (3.9,-6.15) {\texttt{GenerateMissingBases}};

\node[align=center, anchor=west] at (3.05,-7.05) {Public \\ Integration Views:};
\draw[black!80, thick] (3.8, -7.55) -- (3.8, -8.45);
\node[anchor=west] at (3.9,-7.8) {\texttt{ReturnTTerms}};
\node[anchor=west] at (3.9,-8.2) {\texttt{ReturnMasterCombination}};

\node[align=center, anchor=west] at (3.05,-9.05) {Diagnostics};
\draw[black!80, thick] (3.8, -9.55) -- (3.8, -10);
\node[anchor=west] at (3.9,-9.8) {\texttt{DetailedTimingDiagnostics}};

\end{tikzpicture}
\caption{\centering Profile-driven architecture of \textit{AntCalc}'s Integration stage. The principal integration-only options are shown inside the red card.}
\label{fig:antcalc-architecture-intstage}
\end{figure}


\subsection{Integration Profiles and Selection}

The integration-specific profile for each route selects the kind of backend process (and
external package workflow used), and the basis family required for the process. The antennae
are divided by topology, that is, number of final-state particles and number of loops, and all
of those that share the same topology in the current implementation, share the same basis-family
set. Since IBP reduction acts on the kinematic integral structure and is independent of the colour
coefficients, the components of a set are reduced separately. \textit{AntCalc} can therefore
integrate the complete set or a single selected component, avoiding unnecessary reductions.
The exception to this rule is the $A_2^2$ set, in which different sets are used depending on
configuration. A diagram of this route's multi-configuration structure is shown in
Figure~\ref{fig:a22-component-extraction}; comparing it with the upper path of
Figure~\ref{fig:component-extraction} shows how the $A_2^2$ route differs from the
single-family routes used for $A_3^1$ and most others.

The $A_2^1$ antenna is also an exception to the use of integral families: its
loop integrals are reduced directly to scalar one-loop integrals using Passarino-Veltman
reduction, whose known analytic expressions are then substituted and expanded as Laurent
series in $\epsilon$. An overview of these is shown in the following Table~\ref{tab:intModes}.

\begin{table}[h!]
\centering
\renewcommand{\arraystretch}{1.3}
  \begin{tabular}{|c|c|c|} 
   \hline
   \textbf{Default Backend} & \textbf{Bases-Family Set} & \textbf{Current Route} \\ 
   \hline
   Passarino-Veltman / \textit{FeynHelpers} & massless two-parton vertex & $\mathcal A_2^1$ \\ \hline
   IBP / \textit{LiteRed2} & \texttt{X30} & $\mathcal A_3^0$ \\ \hline
   IBP / \textit{LiteRed2} & \texttt{X40} & $\mathcal A_4^0$, $\tilde{\mathcal A}_4^0$, $\mathcal B_4^0$, $\mathcal C_4^0$ \\ \hline
   IBP / \textit{LiteRed2} & \texttt{A31} & $\mathcal A_3^1$, $\tilde{\mathcal A}_3^1$, $\hat{\mathcal A}_3^1$ \\ \hline
   IBP / \textit{LiteRed2} & \texttt{A22TwoLoopTree} & $\mathcal A_2^2$, $\tilde{\mathcal A}_2^2$, $\hat{\mathcal A}_2^2$ \\ \hline
   IBP / \textit{LiteRed2} & \texttt{A22OneLoopSelf} & $\breve{\mathcal A}_2^2$ \\ \hline
   IBP / \textit{LiteRed2} & \texttt{MX30} & $\mathcal A_3^{0,(m_q)}$ \\ \hline
  \end{tabular}
  
  \caption{Overview of the integration-side route-profiles.}
\label{tab:intModes}
\end{table}

\begin{figure}[h!]
\centering
\begin{tikzpicture}[
  font=\small,
  >=Stealth,
  stage/.style={
    draw=black!65,
    rounded corners=2pt,
    fill=white,
    align=center,
    minimum height=10mm,
    text width=19mm,
    inner sep=3pt
  },
  profile/.style={
    stage,
    fill=violet!7,
    text width=34mm,
    minimum height=12mm
  },
  public/.style={
    stage,
    fill=blue!7,
    text width=24mm
  },
  flow/.style={->, line width=.75pt},
  config/.style={->, dashed, gray!75, line width=.6pt}
]

% centre line
\draw[black, thick] (-0.5,-1) -- (1,-1);

\draw[black, thick] (10,-1) -- (11.5,-1);

% build-only lines
\draw[black, thick] (3,-2) -- (6.5,-2);
\draw[black, thick] (7.5,-2) -- (8,-2);

% build-only nodes

% integrate-only lines
\draw[black, thick] (3,0) -- (5,0);

\draw[black, thick] (6,0.5) -- (6.5,0.5);
\draw[black, thick] (6,0) -- (6.5,0);
\draw[black, thick] (6,-0.5) -- (6.5,-0.5);

\draw[black, thick] (7.5,0.5) -- (8,0.5);
\draw[black, thick] (7.5,0) -- (8,0);
\draw[black, thick] (7.5,-0.5) -- (8,-0.5);

% integrate-only nodes
\filldraw[black] (7,0.5) node[anchor=center]{$\mathcal A_2^2$};
\filldraw[black] (7,0) node[anchor=center]{$\tilde{\mathcal A}_2^2$};
\filldraw[black] (7,-0.5) node[anchor=center]{$\hat{\mathcal A}_2^2$};
\filldraw[black] (7,-2) node[anchor=center]{$\breve{\mathcal A}_2^2$};

% curvy lines
\draw (1,-1) to[out=0,in=180] (3,-2);
\draw (1,-1) to[out=0,in=180] (3,0);

\draw (5,0) to[out=0,in=180] (6,0.5);
\draw (5,0) to[out=0,in=180] (6,0);
\draw (5,0) to[out=0,in=180] (6,-0.5);

\draw (8,0.5) to[out=0,in=180] (10,-1);
\draw (8,0) to[out=0,in=180] (10,-1);
\draw (8,-0.5) to[out=0,in=180] (10,-1);
\draw (8,-2) to[out=0,in=180] (10,-1);

% back and frontend backgrounds
\begin{pgfonlayer}{background}
  % Draws a light blue rectangle behind a specific region
  \fill[black!10, rounded corners] (-2.5,1.5) rectangle (13.5,-3.5);
\end{pgfonlayer}

% diagram labels
% \node[label={[text=gray]right:\textit{PaVe}}] at (-2,-4) {};

% circle points
\filldraw[black] (-0.5,-1) circle (2pt) node[align=center, anchor=south, yshift = 2pt]{$\{A_2^2, \tilde A_2^2, \hat A_2^2, \breve A_2^2\}$};
\filldraw[black] (-0.5,-1) circle (2pt) node[align=center, anchor=north, yshift = -2pt]{Key = \texttt{[A,2,2]}};

\filldraw[black] (11.5,-1) circle (2pt) node[align=center, anchor=south, yshift = 2pt]{$\{\mathcal A_2^2, \tilde{\mathcal A}_2^2, \hat{\mathcal A}_2^2, \breve{\mathcal A}_2^2\}$};

% arrows
\node[draw, fill=black, regular polygon, regular polygon sides=5, inner sep=2pt, label={[align=center]above:\texttt{TwoLoopTree}}] at (4,0) {};
\node[draw, fill=black, regular polygon, regular polygon sides=5, inner sep=2pt, label={[align=center]below:\texttt{OneLoopSelf}}] at (4,-2) {};

\end{tikzpicture}
\caption{\centering Diagram showing the component extraction steps and integration-basis families used in the integration-side workflow of the $A_2^2$ antenna set.}
\label{fig:a22-component-extraction}
\end{figure}


As on the build side, the integration profile does not itself perform any calculation; rather,
it specifies the information required to execute the corresponding integration workflow.
Consider, for example, the integration of the $B_4^0$ antenna. From the route key stored in
the \texttt{AntennaObject}, \textit{AntCalc} identifies the appropriate integration profile and
the integral family associated with the $X_4^0$ phase-space topology. The unintegrated antenna
is then mapped onto this integral family and reduced by IBP, using \textit{LiteRed2}, to the
corresponding set of master integrals. Their known analytic expressions are subsequently
substituted to obtain the integrated antenna.

The master-integral basis used in these reductions is discussed in
Chapter~\ref{ch:workedexample}, where selected master integrals are evaluated explicitly and
their role in the analytic integration of antenna functions is illustrated through
representative examples.



\subsection{Integration Routes}

\subsubsection{Passarino-Veltman Route}

The $A_2^1$ antenna constitutes a special low-multiplicity virtual case and is the only
integration route implemented without IBP reduction. As a two-parton one-loop form-factor
contribution, its loop integrals can be treated directly using Passarino-Veltman reduction,
which reduces the tensor loop integrals to scalar one-loop functions. \textit{AntCalc}
evaluates these scalar functions analytically using \textit{Package-X} \cite{PackageX} through
\textit{FeynHelpers} \cite{FeynHelpers}. The required branch and scale conventions are then
applied, followed by the appropriate kinematic normalisation and Laurent expansion in
$\epsilon$, yielding the integrated $\mathcal A_2^1$ antenna. The workflow is summarised
schematically in Fig.~\ref{fig:pave-reduction}.

\begin{figure}[h!]
\centering
\begin{tikzpicture}[
  font=\small,
  >=Stealth,
  stage/.style={
    draw=black!65,
    rounded corners=2pt,
    fill=white,
    align=center,
    minimum height=10mm,
    text width=19mm,
    inner sep=3pt
  },
  profile/.style={
    stage,
    fill=violet!7,
    text width=34mm,
    minimum height=12mm
  },
  public/.style={
    stage,
    fill=blue!7,
    text width=24mm
  },
  flow/.style={->, line width=.75pt},
  config/.style={->, dashed, gray!75, line width=.6pt}
]

% main lines
\draw[black, thick] (0,-1) -- (10,-1);

% back and frontend backgrounds
\begin{pgfonlayer}{background}
  % Draws a light blue rectangle behind a specific region
  \fill[black!10, rounded corners] (-2,0.5) rectangle (12,-2.5);
\end{pgfonlayer}

% back and frontend labels
\node[label={[text=gray]right:\textit{PaVe}}] at (-2,-2) {};

% circle points
\filldraw[black] (0,-1) circle (2pt) node[anchor=south, yshift = 2pt]{$A_2^1$};
\filldraw[black] (2.5,-1) circle (2pt) node[anchor=south, yshift = 2pt]{PaVe};
\filldraw[black] (5,-1) circle (2pt) node[align=center, anchor=south, yshift = 2pt]{Evaluate using \\ \textit{FeynHelpers}};
\filldraw[black] (7.5,-1) circle (2pt) node[align=center, anchor=south, yshift = 2pt]{Laurent- \\ -Expansion};
\filldraw[black] (10,-1) circle (2pt) node[anchor=south, yshift = 2pt]{$\mathcal A_2^1$};

% arrows
\node[draw, fill=black, isosceles triangle, isosceles triangle apex angle=60, inner sep=2pt, shape border rotate=-45] at (1.25,-1) {};
\node[draw, fill=black, isosceles triangle, isosceles triangle apex angle=60, inner sep=2pt, shape border rotate=-45] at (3.75,-1) {};
\node[draw, fill=black, isosceles triangle, isosceles triangle apex angle=60, inner sep=2pt, shape border rotate=-45] at (6.25,-1) {};
\node[draw, fill=black, isosceles triangle, isosceles triangle apex angle=60, inner sep=2pt, shape border rotate=-45] at (8.75,-1) {};

\end{tikzpicture}
\caption{PaVe reduction for the $A_2^1$ antenna integration.}
\label{fig:pave-reduction}
\end{figure}


\subsubsection{IBP Route}

All remaining antenna functions are integrated using IBP reduction, implemented through
\textit{LiteRed2} \cite{Litered}, following the formalism described in
Section~\ref{sec:integration}. The antenna integrand is first expressed in terms of the
momenta, invariants, and denominator structures defining the relevant integral families. The
appropriate phase-space cut propagators are then introduced, allowing the phase-space integrals
to be treated within the same integral-family framework as loop integrals. After expanding the
antenna into individual integral terms, each term is assigned to the corresponding integral
family and reduced by \textit{LiteRed2} to a linear combination of master integrals. The
reduced terms are subsequently recombined to obtain the complete integrated antenna in terms of
the corresponding master-integral basis.

For the public routes, the master integrals (MIs) identified by the IBP reduction are matched
to their corresponding analytic expressions known from the literature and evaluated explicitly
in this work. Their expressions are then substituted by \textit{AntCalc}. Finally, the required
normalisation and dimensional-regularisation conventions are applied, and the result is
expanded as a Laurent series in $\epsilon$, yielding the integrated antenna function. Selected
master integrals and explicit examples of their analytic evaluation are presented in
Chapter~\ref{ch:workedexample}, together with representative applications of this procedure to
the integration of antenna functions. The complete IBP-based integration workflow is summarised
schematically in Fig.~\ref{fig:ibp-reduction}.

\begin{figure}[h!]
\centering
\begin{tikzpicture}[
  font=\small,
  >=Stealth,
  stage/.style={
    draw=black!65,
    rounded corners=2pt,
    fill=white,
    align=center,
    minimum height=10mm,
    text width=19mm,
    inner sep=3pt
  },
  profile/.style={
    stage,
    fill=violet!7,
    text width=34mm,
    minimum height=12mm
  },
  public/.style={
    stage,
    fill=blue!7,
    text width=24mm
  },
  flow/.style={->, line width=.75pt},
  config/.style={->, dashed, gray!75, line width=.6pt}
]

% main lines
\draw[black, thick] (0,-1) -- (10,-1)
                    arc [start angle=90, end angle=-90, radius=0.5]
                           -- (0,-2);

% back and frontend backgrounds
\begin{pgfonlayer}{background}
  % Draws a light blue rectangle behind a specific region
  \fill[black!10, rounded corners] (-2,0.5) rectangle (12.5,-4);
\end{pgfonlayer}

% back and frontend labels
\node[label={[text=gray]right:\textit{IBP}}] at (-2,-3.5) {};

% circle points 
%   top
\filldraw[black] (0,-1) circle (2pt) node[anchor=south, yshift = 2pt]{$X_n^l$};
\filldraw[black] (2.5,-1) circle (2pt) node[align=center, anchor=south, yshift = 2pt]{Mom., inv., den. \\ translation};
\filldraw[black] (5,-1) circle (2pt) node[align=center, anchor=south, yshift = 2pt]{Expansion \\ into terms};
\filldraw[black] (7.5,-1) circle (2pt) node[align=center, anchor=south, yshift = 2pt]{Term \\ matching};
\filldraw[black] (10,-1) circle (2pt) node[align=center, anchor=south, yshift = 2pt]{IBP \\ term-by-term};
%   bottom
\filldraw[black] (10,-2) circle (2pt) node[align=center, anchor=north, yshift = -2pt]{Sum};
\filldraw[black] (7.5,-2) circle (2pt) node[align=center, anchor=north, yshift = -2pt]{Id. \\ and match \\ MI};
\filldraw[black] (5,-2) circle (2pt) node[align=center, anchor=north, yshift = -2pt]{Substitute \\ MI};
\filldraw[black] (2.5,-2) circle (2pt) node[align=center, anchor=north, yshift = -2pt]{Laurent- \\ -Expansion};
\filldraw[black] (0,-2) circle (2pt) node[anchor=north, yshift = -2pt]{$\mathcal X_n^l$};

% arrows
\node[draw, fill=black, isosceles triangle, isosceles triangle apex angle=60, inner sep=2pt, shape border rotate=-45] at (1.25,-1) {};
\node[draw, fill=black, isosceles triangle, isosceles triangle apex angle=60, inner sep=2pt, shape border rotate=-45] at (3.75,-1) {};
\node[draw, fill=black, isosceles triangle, isosceles triangle apex angle=60, inner sep=2pt, shape border rotate=-45] at (6.25,-1) {};
\node[draw, fill=black, isosceles triangle, isosceles triangle apex angle=60, inner sep=2pt, shape border rotate=-45] at (8.75,-1) {};

\node[draw, fill=black, isosceles triangle, isosceles triangle apex angle=60, inner sep=2pt, shape border rotate=-90] at (10.5,-1.5) {};

\node[draw, fill=black, isosceles triangle, isosceles triangle apex angle=60, inner sep=2pt, shape border rotate=-225] at (1.25,-2) {};
\node[draw, fill=black, isosceles triangle, isosceles triangle apex angle=60, inner sep=2pt, shape border rotate=-225] at (3.75,-2) {};
\node[draw, fill=black, isosceles triangle, isosceles triangle apex angle=60, inner sep=2pt, shape border rotate=-225] at (6.25,-2) {};
\node[draw, fill=black, isosceles triangle, isosceles triangle apex angle=60, inner sep=2pt, shape border rotate=-225] at (8.75,-2) {};

% compute MI line
\draw[black, thick] (9.875,-3) -- (10.5,-3);
\draw (9.875,-3) to[out=180,in=0] (7.5,-2);
\filldraw[black] (10.5,-3) circle (2pt) node[anchor=north, yshift = -2pt]{Compute MI};
\node[draw, fill=black, isosceles triangle, isosceles triangle apex angle=60, inner sep=2pt, rotate=-95] at (8.8,-2.55) {};

\end{tikzpicture}
\caption{\centering IBP reduction workflow for the remaining antennae integrations. The MI computation is external to \textit{AntCalc}'s internal operation, and hence is represented by an external incoming arrow.}
\label{fig:ibp-reduction}
\end{figure}


\iffalse

\subsection{Normalisations and Convention Bridge}

As discussed in Section~\ref{sec:dimreg}, before integration, the antenna functions must be 
normalised using the $\mathcal N(\epsilon, k)$ factor, introduced in Eq.~\ref{eq:normN}. 
Since in Section~\ref{sec:buildStage}, where we discussed the build stage, the unintegrated
antenna functions were normalised using the loop-order factor, $\Lambda_l$, in this stage, 
the rest of the normalisations and conventions must be applied. That begins with the 
application of what remains from the $G_k$ factor and then $C(\epsilon, k)$. Finally, we 
reach the full $\mathcal N(\epsilon, k)$ by dividing by the two-particle phase-space measure. 
This is, schematically,
\begin{equation}
    \begin{split}
        \mathcal X_n^l &= \mathcal N(\epsilon, k)\int d\Phi_{X_{ijk}}\, X_n^l\\
                       &= \frac{1}{\Phi_2}\left(8\pi^2\right)^{n-2}S_\epsilon\int d\Phi_{X_{ijk}}\,\left(\Lambda_l X_n^l\right).
    \end{split}
\end{equation}
Note, however, that not all routes require that the rest of the $G_k$ be applied, based on 
their $n$. There are also routes that did not require the application of the $\Lambda_l$
factor in the build stage. The following Table~\ref{tab:GkNorm} shows how each antenna or
antenna set interacts with $G_k$.

\begin{table}[h!]
\centering
\renewcommand{\arraystretch}{1.3}
  \begin{tabular}{|c|c|c|} 
   \hline
   \textbf{Route} & $n$, $l$ & $G_k$\textbf{ relative to build }$\Lambda_l$ \\ 
   \hline
   $A_2^1$ & $2,1$ & $G_1=\Lambda_1$ (no extra multiplicity factors needed) \\ \hline
   $A_3^0$ & $3,0$ & $G_1=8\pi^2\Lambda_0=8\pi^2$ \\ \hline
   $A_3^1$ set & $3,1$ & $G_2=8\pi^2\Lambda_1$ \\ \hline
   $A_4^0$ set, $B_4^0$, $C_4^0$ & $4,0$ & $G_2=(8\pi^2)^2\Lambda_0=(8\pi^2)^2$ \\ \hline
   $A_2^2$ set & $2,2$ & $G_2=\Lambda_2$ (no extra multiplicity factors needed) \\ \hline
  \end{tabular}
  
  \caption{Overview of the interactions of the antenna routes with the normalisation factor $G_k$.}
\label{tab:GkNorm}
\end{table}

\subsubsection{Expansion Order}

The final step in obtaining an integrated antenna $\mathcal X_n^l$, or an ordered list of its
components, is to expand the result as a Laurent series in $\epsilon$. The default truncation
order depends on the perturbative order of the antenna. Following the conventions used in the
literature, the NLO antennae $\mathcal A_2^1$ and $\mathcal A_3^0$ are provided through
$\mathcal O(\epsilon^2)$, whereas the NNLO antennae are provided through $\mathcal O(\epsilon^0)$.
For the currently implemented routes, the available expansion depth is also determined by the
analytic master-integral expressions used by \textit{AntCalc}, some of which have been
implemented and validated only through finite order. The expansion order can nevertheless be
adjusted when the required higher-order coefficients are available, using the
\texttt{ExpansionOrder} option, as,
\begin{equation}\label{eq:expansionOrderDiag}
    \begin{gathered}
        \texttt{IntegrateAntenna[A30]} \xrightarrow{\text{Default Result}}\frac{1}{\epsilon^2}+\frac{3}{2\epsilon}+\frac{19}{4}-\frac{7\pi^2}{12}+\epsilon\!\left(\frac{109}{8}+...\right)+\epsilon^2\!\left(\frac{639}{16}+...\right)\\
        \texttt{IntegrateAntenna[A30]} \xrightarrow{\texttt{ExpansionOrder->0}}\frac{1}{\epsilon^2}+\frac{3}{2\epsilon}+\frac{19}{4}-\frac{7\pi^2}{12}
    \end{gathered}
\end{equation}
where \texttt{A30=BuildAntenna[A,3,0,IntegrableForm -> True]}, is the $A_3^0$ set's \\
\texttt{AntennaObject}. Equation~\ref{eq:expansionOrderDiag} illustrates the use of
\texttt{ExpansionOrder} to override the default truncation order. The option can request
either lower or higher orders; however, the maximum accessible order is route-dependent and is
limited by the available $\epsilon$-expansions of the corresponding master integrals. For the
example shown, all required terms are available and the result has been validated against the
literature.

\subsection{Reconstruction of Colour Components}

For antenna sets containing several colour structures, the corresponding components are
identified during the build stage and kept separate throughout the integration workflow. For
example, the $A_3^1$ set contains leading-colour, subleading-colour, and closed-quark-loop
components, corresponding to $A_3^1$, $\tilde A_3^1$, and $\hat A_3^1$, respectively. Since IBP
reduction acts on the kinematic integral structure and is independent of the colour
coefficients, the components can be treated separately during the reduction. \textit{AntCalc}
can therefore integrate either the complete set of components or only a selected component. The
latter avoids unnecessary IBP reductions and can significantly reduce the computational cost.
Keeping the components separate also allows each contribution to be mapped to the integral
families and master integrals required by its particular kinematic structure.
Fig.~\ref{fig:component-extraction} illustrates this procedure for the $A_3^1$ antenna set.

\begin{figure}[h!]
\centering
\begin{tikzpicture}[
  font=\small,
  >=Stealth,
  stage/.style={
    draw=black!65,
    rounded corners=2pt,
    fill=white,
    align=center,
    minimum height=10mm,
    text width=19mm,
    inner sep=3pt
  },
  profile/.style={
    stage,
    fill=violet!7,
    text width=34mm,
    minimum height=12mm
  },
  public/.style={
    stage,
    fill=blue!7,
    text width=24mm
  },
  flow/.style={->, line width=.75pt},
  config/.style={->, dashed, gray!75, line width=.6pt}
]

% centre line
\draw[black, thick] (0,-1) -- (2,-1);

% build-only lines
\draw[black, thick] (4,-2) -- (5,-2);
\draw[black, dashed] (5,-2) -- (6,-2);

\draw[black, dashed] (7, -1.5) -- (7.5,-1.5);
\draw[black, dashed] (7, -2) -- (7.5,-2);
\draw[black, dashed] (7, -2.5) -- (7.5,-2.5);

% build-only nodes
\filldraw[black] (8,-1.5) node[anchor=center]{$A_3^1$};
\filldraw[black] (8,-2) node[anchor=center]{$\tilde A_3^1$};
\filldraw[black] (8,-2.5) node[anchor=center]{$\hat A_3^1$};

% integrate-only lines
\draw[black, thick] (4,0) -- (6,0);

\draw[black, thick] (7,0.5) -- (7.5,0.5);
\draw[black, thick] (7,0) -- (7.5,0);
\draw[black, thick] (7,-0.5) -- (7.5,-0.5);

\draw[black, dashed] (8.5,0.5) -- (9,0.5);
\draw[black, dashed] (8.5,0) -- (9,0);
\draw[black, dashed] (8.5,-0.5) -- (9,-0.5);

\draw[black, dashed] (10,0) -- (10.5,0)
                    arc [start angle=90, end angle=0, radius=0.5]
                           -- (11,-2);

% integrate-only nodes
\filldraw[black] (8,0.5) node[anchor=center]{$\mathcal A_3^1$};
\filldraw[black] (8,0) node[anchor=center]{$\tilde{\mathcal A}_3^1$};
\filldraw[black] (8,-0.5) node[anchor=center]{$\hat{\mathcal A}_3^1$};

% curvy lines
\draw (2,-1) to[out=0,in=180] (4,-2);
\draw (2,-1) to[out=0,in=180] (4,0);

\draw (6,0) to[out=0,in=180] (7,0.5);
\draw (6,0) to[out=0,in=180] (7,0);
\draw (6,0) to[out=0,in=180] (7,-0.5);

\draw[black, dashed] (6,-2) to[out=0,in=180] (7, -1.5);
\draw[black, dashed] (6,-2) to[out=0,in=180] (7, -2);
\draw[black, dashed] (6,-2) to[out=0,in=180] (7, -2.5);

\draw[black, dashed] (9,0.5) to[out=0,in=180] (10, 0);
\draw[black, dashed] (9,0) to[out=0,in=180] (10, 0);
\draw[black, dashed] (9,-0.5) to[out=0,in=180] (10, 0);

% back and frontend backgrounds
\begin{pgfonlayer}{background}
  % Draws a light blue rectangle behind a specific region
  \fill[black!10, rounded corners] (-2,1.5) rectangle (14,-4);
\end{pgfonlayer}

% diagram labels
% \node[label={[text=gray]right:\textit{PaVe}}] at (-2,-4) {};

% circle points
\filldraw[black] (0,-1) circle (2pt) node[anchor=south, yshift = 2pt]{BuildAntenna[A,3,1]};
\filldraw[black] (11,-2) circle (2pt) node[align=center, anchor=north, yshift = -2pt]{$\{\mathcal A_3^1, \tilde{\mathcal A}_3^1, \hat{\mathcal A}_3^1\}$};
\filldraw[black] (5,-2) circle (2pt) node[anchor=north, yshift = -2pt]{$\{A_3^1, \tilde A_3^1, \hat A_3^1\}$};

% arrows
\node[draw, fill=black, isosceles triangle, isosceles triangle apex angle=60, inner sep=2pt, shape border rotate=-45, label={[align=center]above:\textit{AntennaObject}}] at (5,0) {};
\node[draw, isosceles triangle, isosceles triangle apex angle=60, inner sep=2pt, shape border rotate=-90] at (11,-1) {};

\end{tikzpicture}
\caption{\centering Diagram showing the component extraction steps and paths for the build-only and build-and-integrate routes. The combined route is shown above and the build-only route below it. The dashed lines represent the options to either separate the ordered lists into single components on the build-only route, or to merge the components into a single ordered list for the combined route. The $A_3^1$ antenna set is used as an example.}
\label{fig:component-extraction}
\end{figure}


It is also important to note here that the result of the $A_3^1$ set's 
integration is given, much like it was for the build stage in \ref{eq:a31Norm}, as,
\begin{equation}\label{eq:a31WithTs}
    \begin{split}
        &\mathcal A_3^1 = \mathcal T_\text{Lead} - \mathcal A_2^1\mathcal A_3^0\\
        &\tilde{\mathcal A}_3^1 = -\left(\mathcal T_\text{Sublead} + \mathcal A_2^1\mathcal A_3^0\right)\\
        &\hat{\mathcal A}_3^1 = \mathcal T_\text{QL}.
    \end{split}
\end{equation}
These resulting $\mathcal A_3^1$ and $\tilde{\mathcal A}_3^1$ are already
UV renormalised by this stage.

Those $\mathcal T$-object may still be useful, however. They can be extracted instead of the 
antenna functions, essentially by breaking the process before that last operation, in 
Eq.~\ref{eq:a31WithTs} and returning an analogous ordered list made up of the $\mathcal T$-objects,
as,
\begin{equation}
    \begin{gathered}
        \texttt{IntegrateAntenna[A31]} \xrightarrow{\text{Default Result}} \{\mathcal A_3^1, \tilde{\mathcal A}_3^1, \hat{\mathcal A}_3^1\}\\
        \texttt{IntegrateAntenna[A31]} \xrightarrow{\texttt{ReturnTTerms -> True}} \{\mathcal T_\text{Lead}, \mathcal T_\text{Sublead}, \mathcal T_\text{QL}\}
    \end{gathered}
\end{equation}
where \texttt{A31=BuildAntenna[A,3,1,IntegrableForm -> True]}, is, once again, the $A_3^1$ set's
\texttt{AntennaObject}.

\subsection{$A_2^2$ Set Integration Routes}

As summarised in Table~\ref{tab:twoLoopProdModes}, the $A_2^2$ antenna set is divided into two
subsets, and require distinct integral families for the IBP reduction. Consequently, the two
subsets follow separate integration routes in \textit{AntCalc}.
Those routes may be schematically drawn in the spirit of the diagram in
Fig.~\ref{fig:component-extraction}, as in the diagram in 
Fig.~\ref{fig:a22-component-extraction}, below. The workflow interprets the overall 
\texttt{[A,2,2]}, then separates the process into different build-side production modes 
as shown in the diagram in Fig.~\ref{fig:a22-build}. It then merges the four components into 
an ordered list. That list is sent to the integrate-side workflow. This workflow starts by 
separating them into the subsets and then, considering that all components are selected for 
integration, it reduces them using the subset-specific basis families. Finally, once all four 
are integrated, they are once again merged onto an ordered list, which is delivered as the 
default output.

\begin{figure}[h!]
\centering
\begin{tikzpicture}[
  font=\small,
  >=Stealth,
  stage/.style={
    draw=black!65,
    rounded corners=2pt,
    fill=white,
    align=center,
    minimum height=10mm,
    text width=19mm,
    inner sep=3pt
  },
  profile/.style={
    stage,
    fill=violet!7,
    text width=34mm,
    minimum height=12mm
  },
  public/.style={
    stage,
    fill=blue!7,
    text width=24mm
  },
  flow/.style={->, line width=.75pt},
  config/.style={->, dashed, gray!75, line width=.6pt}
]

% centre line
\draw[black, thick] (-0.5,-1) -- (1,-1);

\draw[black, thick] (10,-1) -- (11.5,-1);

% build-only lines
\draw[black, thick] (3,-2) -- (6.5,-2);
\draw[black, thick] (7.5,-2) -- (8,-2);

% build-only nodes

% integrate-only lines
\draw[black, thick] (3,0) -- (5,0);

\draw[black, thick] (6,0.5) -- (6.5,0.5);
\draw[black, thick] (6,0) -- (6.5,0);
\draw[black, thick] (6,-0.5) -- (6.5,-0.5);

\draw[black, thick] (7.5,0.5) -- (8,0.5);
\draw[black, thick] (7.5,0) -- (8,0);
\draw[black, thick] (7.5,-0.5) -- (8,-0.5);

% integrate-only nodes
\filldraw[black] (7,0.5) node[anchor=center]{$\mathcal A_2^2$};
\filldraw[black] (7,0) node[anchor=center]{$\tilde{\mathcal A}_2^2$};
\filldraw[black] (7,-0.5) node[anchor=center]{$\hat{\mathcal A}_2^2$};
\filldraw[black] (7,-2) node[anchor=center]{$\breve{\mathcal A}_2^2$};

% curvy lines
\draw (1,-1) to[out=0,in=180] (3,-2);
\draw (1,-1) to[out=0,in=180] (3,0);

\draw (5,0) to[out=0,in=180] (6,0.5);
\draw (5,0) to[out=0,in=180] (6,0);
\draw (5,0) to[out=0,in=180] (6,-0.5);

\draw (8,0.5) to[out=0,in=180] (10,-1);
\draw (8,0) to[out=0,in=180] (10,-1);
\draw (8,-0.5) to[out=0,in=180] (10,-1);
\draw (8,-2) to[out=0,in=180] (10,-1);

% back and frontend backgrounds
\begin{pgfonlayer}{background}
  % Draws a light blue rectangle behind a specific region
  \fill[black!10, rounded corners] (-2.5,1.5) rectangle (13.5,-3.5);
\end{pgfonlayer}

% diagram labels
% \node[label={[text=gray]right:\textit{PaVe}}] at (-2,-4) {};

% circle points
\filldraw[black] (-0.5,-1) circle (2pt) node[align=center, anchor=south, yshift = 2pt]{$\{A_2^2, \tilde A_2^2, \hat A_2^2, \breve A_2^2\}$};
\filldraw[black] (-0.5,-1) circle (2pt) node[align=center, anchor=north, yshift = -2pt]{Key = \texttt{[A,2,2]}};

\filldraw[black] (11.5,-1) circle (2pt) node[align=center, anchor=south, yshift = 2pt]{$\{\mathcal A_2^2, \tilde{\mathcal A}_2^2, \hat{\mathcal A}_2^2, \breve{\mathcal A}_2^2\}$};

% arrows
\node[draw, fill=black, regular polygon, regular polygon sides=5, inner sep=2pt, label={[align=center]above:\texttt{TwoLoopTree}}] at (4,0) {};
\node[draw, fill=black, regular polygon, regular polygon sides=5, inner sep=2pt, label={[align=center]below:\texttt{OneLoopSelf}}] at (4,-2) {};

\end{tikzpicture}
\caption{\centering Diagram showing the component extraction steps and integration-basis families used in the integration-side workflow of the $A_2^2$ antenna set.}
\label{fig:a22-component-extraction}
\end{figure}


\fi

\section{Higher-Level Functions}

Besides the two canonical functions, \texttt{BuildAntenna} and \texttt{IntegrateAntenna},
\textit{AntCalc} also supports higher-level functions. This layer does not
introduce new antenna-construction or integration methodology; instead it orchestrates the
two canonical functions to compute observables such as the $R$-ratio, to assemble the
topology-grouped combinations used in it, and to streamline repeated build or integrate
evaluation. The five higher-level functions are:

\begin{itemize}
    \item \texttt{BuildAndIntegrateAntenna} --- chains \texttt{BuildAntenna[\ldots,
    IntegrableForm -> True]} into \texttt{IntegrateAntenna} in a single call, forwarding the
    component-selection and Laurent-expansion-order controls to their respective stages.
    \item \texttt{BuildAllAntennae} --- builds every public antenna function in one call.
    \item \texttt{BuildAndIntegrateAllAntennae} --- builds and integrates every public antenna
    function in one call. Both bulk wrappers can be limited to a chosen perturbative order.
    \item \texttt{TObject} --- assembles integrated antennae into the topology-grouped
    combinations used in the $R$-ratio, defined in Subsection~\ref{subsec:Tobjects} of
    Chapter~\ref{ch:validation}.
    \item \texttt{BuildRRatio} --- assembles the massless hadronic $R$-ratio from the
    integrated antenna ingredients required at the selected perturbative order.
\end{itemize}

The $R$-ratio produced by \texttt{BuildRRatio} is the observable used for \textit{AntCalc}'s
global validation. Called as
\texttt{BuildRRatio[SMQCD\footnote{The label SMQCD is not standard in the literature. It is
used here to denote the massless QCD sector of the Standard Model.}, maxOrder -> \ldots]}, it
generates the $R$-ratio through NNLO and, by default, returns the finite coefficient after
application of the observable convention bridge. This result has been validated against a
literature reference target, as discussed in Chapter~\ref{ch:validation}.

\iffalse
% Parked pending review --- higher-level-functions schematic. With the per-function prose
% walkthroughs removed, this figure illustrates only self-evident orchestration; kept here in
% case a trimmed version is wanted.
\begin{figure}[h!]
\centering
\begin{tikzpicture}[
  font=\small,
  >=Stealth,
  stage/.style={
    draw=black!65,
    rounded corners=2pt,
    fill=white,
    align=center,
    minimum height=10mm,
    text width=19mm,
    inner sep=3pt
  },
  profile/.style={
    stage,
    fill=violet!7,
    text width=34mm,
    minimum height=12mm
  },
  public/.style={
    stage,
    fill=blue!7,
    text width=24mm
  },
  flow/.style={->, line width=.75pt},
  config/.style={->, dashed, gray!75, line width=.6pt}
]

% buildandintegrateantenna
\filldraw[black] (3,0) circle (2pt) node[anchor=south, yshift = 2pt]{BuildAndIntegrateAntenna};

% buildantenna
\node[draw, fill=black, isosceles triangle, isosceles triangle apex angle=60, inner sep=2pt, rotate=-90] at (3,-0.75) {};
\draw[black, thick] (3,0) -- (3,-1.5);
\filldraw[black] (3,-1.5) circle (2pt) node[anchor=west, xshift=2pt]{BuildAntenna};

% build secondary path
\draw[black, thick] (3,-1.5) to[out=-90, in=20] (2,-2.5);
\node[draw, fill=black, isosceles triangle, isosceles triangle apex angle=60, inner sep=2pt, rotate=200, label={[anchor=north east, xshift=-5pt, yshift=9pt]:$X_n^l$}] at (2,-2.5) {};

% integrateantenna
\node[draw, fill=black, isosceles triangle, isosceles triangle apex angle=60, inner sep=2pt, rotate=-90, label={[anchor=west,xshift=2pt]:\textit{AntennaObject}}] at (3,-2.5) {};
\draw[black, thick] (3,-1.5) -- (3,-3.5);
\filldraw[black] (3,-3.5) circle (2pt) node[anchor=west, xshift=2pt]{IntegrateAntenna};

% integrate paths
\draw[black, thick] (3,-3.5) -- (3,-4);
\draw[black, thick] (3,-4) to[out=-90,in=90] (4,-5);
\draw[black, thick] (3,-4) -- (3, -5);
\draw[black, thick] (3,-4) to[out=-90,in=90] (2,-5);
\filldraw[black] (2,-5.5) node[anchor=center]{IBP};
\filldraw[black] (3,-5.5) node[anchor=center]{PaVe};
\filldraw[black] (4,-5.5) node[align=center, anchor=center]{Trivial \\ ($A_2^0$)};
\draw[black, thick] (2,-6) to[out=-90,in=90] (3,-7);
\draw[black, thick] (3,-6) -- (3, -7);
\draw[black, thick] (4,-6) to[out=-90,in=90] (3,-7);
\draw[black, thick] (3,-7) -- (3,-7.5);

% integrate sec path
\draw[black, thick] (3,-7) to[out=-90, in=20] (2,-8);
\node[draw, fill=black, isosceles triangle, isosceles triangle apex angle=60, inner sep=2pt, rotate=200, label={[anchor=north east, xshift=-5pt, yshift=9pt]:$\mathcal X_n^l$}] at (2,-8) {};

% buildandintegrateallantennae
\draw[black, thick] (3,-7.5) to[out=-90, in=90] (5,-10.5);

\filldraw[black] (8,-8.1) node[anchor=center]{$\mathcal C_4^0$};
\filldraw[black] (7.4,-8.1) node[anchor=center]{$\mathcal B_4^0$};
\filldraw[black] (6.8,-8.1) node[anchor=center]{\ldots};
\filldraw[black] (6.2,-8.1) node[anchor=center]{$\mathcal A_2^1$};
\filldraw[black] (5.6,-8.1) node[anchor=center]{$\mathcal A_2^0$};
\draw[black, thick] (8,-8.6) to[out=-90, in=90] (5,-10.5);
\draw[black, thick] (7.4,-8.6) to[out=-90, in=90] (5,-10.5);
\draw[black, thick] (6.8,-8.6) to[out=-90, in=90] (5,-10.5);
\draw[black, thick] (6.2,-8.6) to[out=-90, in=90] (5,-10.5);
\draw[black, thick] (5.6,-8.6) to[out=-90, in=90] (5,-10.5);

% tobject
\draw[black, thick] (3,-7.5) to[out=-90, in=90] (1,-10.5);
\filldraw[black] (1,-10.5) circle (2pt) node[anchor=north, yshift=-2pt]{TObject};

\filldraw[black] (0,-8.1) node[anchor=center]{$\tilde{\mathcal X}_n^l$};
\filldraw[black] (-0.8,-8.1) node[anchor=center]{$\hat{\mathcal X}_n^l$};
\filldraw[black] (-1.6,-8.1) node[anchor=center]{\ldots};
\draw[black, thick] (0,-8.6) to[out=-90, in=90] (1,-10.25);
\draw[black, thick] (-0.8,-8.6) to[out=-90, in=90] (1,-10.25);
\draw[black, thick] (-1.6,-8.6) to[out=-90, in=90] (1,-10.25);

% RRatio
\draw[black, thick] (5,-10.5) to[out=-90, in=90] (8,-12.5);
\filldraw[black] (8,-12.5) circle (2pt) node[align=center, anchor=north, yshift=-2pt]{BuildRRatio};
\draw[black, thick] (5,-10.5) to[out=-90, in=90] (2,-12.5);
\filldraw[black] (2,-12.5) circle (2pt) node[align=center, anchor=north, yshift=-2pt]{BuildAndIntegrateAllAntennae \\ $\{\mathcal A_2^0, \mathcal A_2^1, ..., \mathcal X_n^l, ..., \mathcal B_4^0, \mathcal C_4^0\}$};

% buildallantenna
\filldraw[black] (10,-1.5) circle (2pt) node[anchor=south, yshift = 2pt]{BuildAllAntennae};
\draw[black, thick] (10,-1.5) to[out=-90, in=90] (8.5,-3);
\draw[black, thick] (10,-1.5) to[out=-90, in=90] (9.5,-3);
\draw[black, thick] (10,-1.5) to[out=-90, in=90] (10.5,-3);
\draw[black, thick] (10,-1.5) to[out=-90, in=90] (11.5,-3);
\filldraw[black] (8.5,-3.5) node[anchor=center]{$A_2^0$};
\filldraw[black] (9.5,-3.5) node[anchor=center]{$A_2^1$};
\filldraw[black] (10.5,-3.5) node[anchor=center]{\ldots};
\filldraw[black] (11.5,-3.5) node[anchor=center]{$C_4^0$};
\draw[black, thick] (8.5,-4) to[out=-90, in=90] (10,-5.5);
\draw[black, thick] (9.5,-4) to[out=-90, in=90] (10,-5.5);
\draw[black, thick] (10.5,-4) to[out=-90, in=90] (10,-5.5);
\draw[black, thick] (11.5,-4) to[out=-90, in=90] (10,-5.5);
\filldraw[black] (10,-5.5) circle (2pt) node[anchor=north, yshift = -2pt]{$\{A_2^0, A_2^1,...,B_4^0,C_4^0\}$};

% curvy lines
%\draw (1,-1) to[out=0,in=180] (3,-2);
%\draw (1,-1) to[out=0,in=180] (3,0);

% back and frontend backgrounds
\begin{pgfonlayer}{background}
  % Draws a light blue rectangle behind a specific region
  \fill[black!10, rounded corners] (-2.5,1) rectangle (13.5,-14);
\end{pgfonlayer}

% diagram labels
% \node[label={[text=gray]right:\textit{PaVe}}] at (-2,-4) {};

% circle points

\end{tikzpicture}
\caption{\centering Diagram showing the build-side workflow leading to the generation of the $A_2^2$ antenna set. The two subsets are separated by the interferences used in their generation.}
\label{fig:higher-level}
\end{figure}

\fi

\iffalse
% Parked pending review --- the TObject table and worked example. T-objects are defined in
% Chapter 5 (subsec:Tobjects); this material may belong there rather than in Chapter 3.
We may also retrieve a specific $\mathcal T$-object directly, using \texttt{TObject}. The public supported
$\mathcal T$-objects are described in Table~\ref{tab:TObject}, below. These computation routes follow closely the
methods discussed in Subsection~\ref{subsec:Tobjects} of Chapter~\ref{ch:validation}. The function calls the one-shot integrated routes, using \texttt{BuildAndIntegrateAntenna}
for the necessary ingredient antennae, then assembles them into the $\mathcal T$-object
structure.

\begin{table}[h!]
\centering
\renewcommand{\arraystretch}{1.3}
  \begin{tabular}{|c|c|c|} 
   \hline
   \textbf{$\mathcal T$-Object} & \textbf{Ingredients} & \textbf{Call} \\ 
   \hline
   $\mathcal T_{q\bar q}^2$ & Born contribution & \texttt{TObject[2,"qqbar"]} \\ \hline
   $\mathcal T_{q\bar q}^4$ & $\mathcal A_2^1$ & \texttt{TObject[4,"qqbar"]} \\ \hline
   $\mathcal T_{q\bar qg}^4$ & $\mathcal A_3^0$ & \texttt{TObject[4,"qqbarg"]} \\ \hline
   $\mathcal T_{q\bar q}^6$ & $\mathcal A_2^2$, $\tilde {\mathcal A}_2^2$, $\hat {\mathcal A}_2^2$, $\breve {\mathcal A}_2^2$ & \texttt{TObject[6,"qqbar"]} \\ \hline
   $\mathcal T_{q\bar qg}^6$ & $\mathcal A_2^1$, $\mathcal A_3^0$, $\mathcal A_3^1$, $\tilde {\mathcal A}_3^1$, $\hat {\mathcal A}_3^1$ & \texttt{TObject[6,"qqbarg"]} \\ \hline
   $\mathcal T_{q\bar qq^\prime\bar q^\prime}^6$ & $\mathcal B_4^0$ & \texttt{TObject[6,"qqbarqprimeqprimebar"]} \\ \hline
   $\mathcal T_{q\bar qq\bar q}^6$ & $\mathcal B_4^0$, $\mathcal C_4^0$ & \texttt{TObject[6,"qqbarqqbar"]} \\ \hline
   $\mathcal T_{q\bar qgg}^6$ & $\mathcal A_4^0$, $\tilde{\mathcal A}_4^0$ & \texttt{TObject[6,"qqbargg"]} \\ \hline
  \end{tabular}
   \caption{Overview of the TObject routes.}
\label{tab:TObject}
\end{table}
As an example, let us consider $\mathcal T_{q\bar qg}^6$. As shown in Eq.~\ref{eq:a31Norm}
and Table~\ref{tab:TObject}, this $\mathcal T$-object requires the $\mathcal A_3^1$ set, 
$\mathcal A_2^1$, and $\mathcal A_3^0$. Upon calling \texttt{BuildAndIntegrateAntenna}
for those antennae, \texttt{TObject} assembles them as follows,
\begin{equation}
    \mathcal T_{q\bar qg}^6 = \left(N-\frac1N\right)T_{q\bar q}^2\left[N\!\left(\mathcal A_3^1+\mathcal A_2^1\mathcal A_3^0\right)-\frac1N\!\left(\tilde{\mathcal A}_3^1+\mathcal A_2^1\mathcal A_3^0\right)+N_f\hat{\mathcal A}_3^1\right].
\end{equation}
\fi

\section{Validation and Development Diagnostics}\label{sec:validationDiagnostics}

As outlined in the introduction, the implementation is validated at two levels. Locally, each
integrated antenna function is compared against its published Laurent expansion
\cite{Gehrmann-DeRidder:2005btv,Gehrmann-DeRidder:2004ttg} and, where applicable, against the
expected eikonal factors and Altarelli--Parisi splitting functions in the relevant unresolved
limits (Chapter~\ref{ch:workedexample}). Globally, the integrated antennae are assembled into
the inclusive $R$-ratio and compared with the known finite result
(Chapter~\ref{ch:validation}).

As discussed in Section~\ref{sec:profile}, the route system can retain selected diagnostic
information from an evaluation. That makes it easy to locate any inconsistency during a route
calculation, which can then be inspected using the built-in diagnostics tool --- for example,
by retrieving the antenna ingredients of an $R$-ratio order and checking how their
$\epsilon$-poles cancel.

\iffalse
% Task O --- this worked NLO pole-cancellation example moves to Chapter~\ref{ch:validation},
% where it should be recast as the compact pole-cancellation table (by perturbative order and
% NNLO colour/flavour channel) called for in the review, alongside the R-ratio assembly.
% The "release-level observable validation" sentence below overlaps the release-acceptance
% paragraph in Subsection~\ref{subsec:wardAndRelease} and can be dropped when relocating.
One example of its use could be checking how the poles of
each R-ratio order's antennae cancel each other to make a finite result. As an
example, let us consider the SMQCD R-ratio up to NLO. We can run it and then,
in the diagnostics association, see what antennae it uses as ingredients. We may
then simply compare those against one another (in this case there are only two,
$\mathcal A_2^1$ and $\mathcal A_3^0$), and see how their poles cancel. The
check may go as follows. In the following example, the common colour and perturbative
prefactors are omitted, since they do not affect the cancellation of the poles in
$\mathcal A_2^1 + \mathcal A_3^0$.

\begin{mdframed}[backgroundcolor=gray!10, linewidth=0pt]
    \texttt{In[1]:= <<AntCalc`}\\
    \hfill
    \texttt{In[2]:= diagNLORRatio = BuildRRatio[SMQCD,maxOrder -> NLO,\\ReturnDiagnostics -> True][[2]];}\\
    \hfill
    \texttt{In[3]:= ingredients = diagNLORRatio["Ingredients"];}\\
    \hfill
    \texttt{In[4]:= A21 = ingredients["intA21"]; A30 = ingredients["intA30"];}\\
    \hfill
    \texttt{In[5]:= resList = \{\};}\\
    \hfill
    \texttt{In[6]:= \\Do[\\
                        coefA21 = Coefficient[A21,Epsilon,i];\\
                        coefA30 = Coefficient[A30,Epsilon,i];\\
                        AppendTo[resList, coefA21+coefA30]\\
                        ,\\
                        \{i,-2,2\}\\
                        ];}\\
    \hfill
    \texttt{In[7]:= resList}\\
    \hfill
    \texttt{Out[7]:=}
        \begin{equation*}
            \left\{0,0,\frac34,\frac{45}8-6\zeta_3,\frac{383}{16}-\frac{7\pi^2}{16}-\frac{\pi^4}{10}-9\zeta_3\right\}
        \end{equation*}
\end{mdframed}

The first two entries are the coefficients of $\epsilon^{-2}$ and $\epsilon^{-1}$, 
respectively, and vanish as required. The remaining entries are the finite and 
positive-power coefficients, which do not cancel. This diagnostic 
inspection makes the origin of the NLO pole cancellation explicit. The release-level
observable validation additionally compares the raw NNLO $R$-ratio series with the
known pole-free and finite result.
\fi

The \texttt{Ingredients} association is only one part of the diagnostic information. For a
complete retained history of a calculation, \textit{AntCalc} can instead return an
\texttt{AntennaRunRecord}, which records the relevant route stages and metadata. This
information was used to trace the source of discrepancies during development; the resulting
validation comparisons are presented in Chapter~\ref{ch:validation}.

\subsection{Amplitude-Level and Release-Acceptance Validation}\label{subsec:wardAndRelease}

The checks described above validate an antenna function against an external reference: a
published Laurent series, or a known finite observable. \textit{AntCalc} also implements a
complementary validation family that requires no such reference. The function
\texttt{VerifyWardIdentity[type, numFinalParticles, loopOrder]} acts directly on the unsquared
amplitude produced by the build stage, before spin summation, colour decomposition, or
phase-space integration is performed. For a chosen external gluon leg, it replaces that gluon's
polarisation vector by its own momentum, $\varepsilon^\mu(k)\to k^\mu$, while keeping every other
external gluon polarisation explicitly transverse, and checks whether the resulting expression
vanishes. This is the standard amplitude-level Ward identity, applied leg by leg to the raw
generated amplitude rather than inferred from a property of the extracted or integrated antenna.
It is currently applied to every public route with an external gluon leg, that is, to $A_3^0$,
$A_4^0$ and the $A_3^1$ set, and returns a structured pass, fail, or time-out report for each
checked leg. Because it depends only on the internal gauge structure of the generated amplitude,
and not on comparison with any external result, it certifies a property the literature-comparison
checks above cannot: that the amplitude's diagrams were assembled with the correct relative signs
and couplings before extraction even begins. In principle, this makes it applicable to a route for
which no literature result exists at all.

Separately from this route-level testing, \textit{AntCalc}'s development also relies on a
release-acceptance practice. Before a version is prepared, every public route is re-evaluated in a
fresh kernel, with any stored or cached intermediate result disabled, and its build, integration,
and one-shot outputs are compared against each other and against the literature target discussed
above. Each run is recorded, together with its timing and the specific checks it satisfied, as
dated evidence. This practice has already caught, and allowed the correction of, release-preparation
issues that were not visible from the final Laurent-series outputs alone.

\subsection{Runtime and Memory Characteristics}

To characterise the computational cost and performance of the implemented public 
routes, each antenna was evaluated in a fresh \textit{Wolfram Mathematica} kernel. 
Package loading was excluded from the measurements. Both the build 
and integration stages were timed using the \texttt{AbsoluteTiming} function. 
Memory usage was measured using \texttt{MaxMemoryUsed[]}, relative to the 
initial kernel baseline. These measurements were taken using a MacBook Pro (Mac16.8), 
with an Apple M4 Pro 12 core (8 performance and 4 efficiency) CPU, and 24~GB unified 
RAM, running macOS 26.6.2 (build 25G83), \textit{Wolfram Mathematica} 15.0.0 for Mac 
OS X ARM (64-bit), and AntCalc 0.3.0$\beta$2.

The lower-order routes, up to NLO, complete both stages in approximately one second,
with the cost increasing substantially for the NNLO routes, in particular $A_2^2$, 
which surpasses 31 minutes. The $A_4^0$ route has the largest memory requirement, 
reaching 567 MB above the kernel benchmark. The complete results are shown in 
Table~\ref{tab:runtimeBenchmarks}.

\begin{table}[h!]
    \centering
    \small
    \renewcommand{\arraystretch}{1.15}
    \begin{tabular}{|c|r|r|r|r|}
        \hline
        \textbf{Route} &
        \textbf{Build} &
        \textbf{Integration} &
        \textbf{End-to-end} &
        \textbf{Peak memory} \\
        &
        \textbf{time} &
        \textbf{time} &
        \textbf{time} &
        \textbf{above baseline} \\
        \hline
        $A_2^0$ & $0.09\,\mathrm{s}$ & $0.49\,\mathrm{s}$ & $0.58\,\mathrm{s}$ & $22.7\,\mathrm{MiB}$ \\
        $A_3^0$ & $0.30\,\mathrm{s}$ & $0.85\,\mathrm{s}$ & $1.15\,\mathrm{s}$ & $24.8\,\mathrm{MiB}$ \\
        $A_2^1$ & $0.57\,\mathrm{s}$ & $0.40\,\mathrm{s}$ & $0.98\,\mathrm{s}$ & $19.3\,\mathrm{MiB}$ \\
        $A_3^1$ & $74.67\,\mathrm{s}$ & $15\,\mathrm{min}\,27\,\mathrm{s}$ & $16\,\mathrm{min}\,42\,\mathrm{s}$ & $195.3\,\mathrm{MiB}$ \\
        $A_2^2$ & $21.79\,\mathrm{s}$ & $31\,\mathrm{min}\,06\,\mathrm{s}$ & $31\,\mathrm{min}\,27\,\mathrm{s}$ & $227.7\,\mathrm{MiB}$ \\
        $A_4^0$ & $3\,\mathrm{min}\,03\,\mathrm{s}$ & $13\,\mathrm{min}\,51\,\mathrm{s}$ & $16\,\mathrm{min}\,53\,\mathrm{s}$ & $567.3\,\mathrm{MiB}$ \\
        $B_4^0$ & $8.65\,\mathrm{s}$ & $17.27\,\mathrm{s}$ & $25.92\,\mathrm{s}$ & $143.2\,\mathrm{MiB}$ \\
        $C_4^0$ & $5.56\,\mathrm{s}$ & $5\,\mathrm{min}\,26\,\mathrm{s}$ & $5\,\mathrm{min}\,32\,\mathrm{s}$ & $174.5\,\mathrm{MiB}$ \\
        $A_3^{0,(m_q)}$ & $2.81\,\mathrm{s}$ & $0.25\,\mathrm{s}$ & $3.09\,\mathrm{s}$ & $12.07\,\mathrm{MiB}$ \\
        \hline
    \end{tabular}
    \caption{\centering Representative resource measurements for the implemented public
    antenna routes. The build time measures
    \texttt{BuildAntenna[..., IntegrableForm -> True]}, followed by
    \texttt{IntegrateAntenna[...]} for the integration time. Each route was
    evaluated in a fresh Wolfram kernel. Peak memory is the maximum kernel-memory
    increase relative to the initial baseline.}
    \label{tab:runtimeBenchmarks}
\end{table}